\documentclass[12pt]{article}

\usepackage{sbc-template}
\usepackage[brazil]{babel}
\usepackage[utf8]{inputenc}
\usepackage[T1]{fontenc}
\usepackage{url,amsmath,booktabs}
\usepackage[hidelinks]{hyperref}

\newcommand{\nomeautor}{Cledson da Silva Araújo}

\sloppy

\title{Proposta de Projeto de Aplicação\\
Predição de mortalidade neonatal no estado do Rio de Janeiro\\
com dados abertos do SINASC e do SIM}

\author{\nomeautor\inst{1}}

\address{Instituto de Computação -- Universidade Federal Fluminense (UFF)\\
  Niterói -- RJ -- Brasil\\[2mm]
  Disciplina: Introdução ao Aprendizado de Máquina para a Saúde (2026.2)\\
  Prof. Flávio Luiz Seixas}

\begin{document}

\maketitle

% ---------------------------------------------------------------------------
\section{Tema e motivação}

A mortalidade neonatal (óbito de 0 a 27 dias de vida) responde pela maior parte da mortalidade infantil no Brasil e concentra-se em causas evitáveis ligadas à assistência na gestação e no parto. Um escore de risco calculado automaticamente a partir da Declaração de Nascido Vivo (DN) pode priorizar leitos de UTI neonatal, visitas domiciliares na primeira semana e a regulação de gestantes de alto risco.

\section{Artigo de referência}

\textbf{Batista et al. (2021)} \cite{batista2021}, \textit{Neonatal mortality prediction with routinely collected data: a machine learning approach}, BMC Pediatrics 21:322.

Os autores fizeram o linkage probabilístico entre o SINASC e o SIM do município de São Paulo, com dados identificados cedidos pela prefeitura. A coorte tinha 1.128.162 nascidos vivos em 2012--2017 e 7.282 óbitos neonatais. As 23 variáveis da DN foram usadas para treinar cinco algoritmos, com treino em 2012--2016 e teste em 2017. O melhor foi o XGBoost, com AUROC 0,971, AUPRC 0,586 e F1 0,548. Os 5\% de maior risco concentraram mais de 90\% dos óbitos, e as variáveis mais importantes foram Apgar, peso, anomalia congênita e idade gestacional.

\textbf{Lacunas.}
\begin{itemize}
  \item O estudo usa uma única cidade e dados identificados, o que impede a reprodução por terceiros.
  \item O modelo depende de variáveis que só existem depois do parto (Apgar, peso).
  \item Não avalia calibração nem desempenho por subgrupo. Essas duas lacunas também são apontadas pela revisão sistemática de \cite{mangold2021}: só 45\% dos estudos reportam calibração e 18\% fazem validação externa.
\end{itemize}

\section{Perguntas de pesquisa e hipóteses}

\begin{enumerate}
  \item \textbf{Reprodução com dados abertos.} O desempenho do artigo se mantém no estado do RJ usando apenas microdados públicos, sem nenhum identificador? \textit{H1: AUROC $\geq$ 0,95 no modelo completo.}
  \item \textbf{Risco antes do parto.} Quanto se perde com um modelo que usa só o que se sabe na admissão para o parto (mãe, gestação e pré-natal, sem Apgar, peso nem anomalia)? \textit{H2: queda substancial de AUROC, com o modelo ainda útil para triagem.}
  \item \textbf{Calibração e equidade.} O desempenho e a calibração variam por raça/cor e escolaridade da mãe e entre capital e interior? \textit{H3: a calibração piora nos subgrupos de menor escolaridade.}
\end{enumerate}

\section{Dados}

\begin{itemize}
  \item \textbf{Fontes:} microdados públicos do SINASC (nascidos vivos) e do SIM (óbitos), por residência no RJ, obtidos no FTP do DATASUS sem cadastro \cite{datasus}. O download é automatizado.
  \item \textbf{Período:} 1.560.938 nascidos vivos em 2017--2024 e óbitos de 2017--2025. O ano de 2025 (SIM ainda preliminar) entra só para captar óbitos de nascidos em dezembro de 2024.
  \item \textbf{Desfecho:} óbito não fetal com 0--27 dias de vida.
\end{itemize}

\textbf{Linkage sem identificadores.} Os arquivos públicos não trazem o número da DN. Porém, a Declaração de Óbito de menores de 1 ano repete dados do nascimento. Cada óbito é pareado a um nascido vivo por concordância exata em seis chaves compostas em cascata:
\begin{itemize}
  \item a mais restritiva é data de nascimento + sexo + município + peso + idade da mãe + semanas de gestação + tipo de gravidez + tipo de parto;
  \item as seguintes toleram campos ausentes;
  \item em todas, só se aceitam pares 1:1.
\end{itemize}

Os resultados do linkage já foram verificados:
\begin{itemize}
  \item 12.155 dos 13.471 óbitos neonatais do período foram pareados (\textbf{90,2\%}, entre 88,5\% e 91,6\% conforme o ano).
  \item Esse percentual é comparável aos 90,8\% obtidos com dados identificados nas capitais brasileiras \cite{maia2017}.
  \item A fração de falsos pares, estimada por um modelo de mistura sobre campos fora das chaves, é de 6--8\% nas duas chaves principais.
\end{itemize}
Essa estimativa usa a escolaridade da mãe e as perdas fetais anteriores, tomando como referência a passada completa e pares aleatórios.

\section{Método}

\begin{itemize}
  \item \textbf{Conjuntos de variáveis:}
    \begin{itemize}
      \item (A) \textit{completo}: as 28 variáveis da DN, como no artigo;
      \item (B) \textit{pré-parto}: 17 variáveis da mãe, da gestação e do pré-natal (idade gestacional, consultas, índice de Kotelchuck, histórico obstétrico).
    \end{itemize}
  \item \textbf{Modelos:} regressão logística (linha de base interpretável), XGBoost (modelo do artigo) \cite{chen2016} e LightGBM \cite{ke2017}.
  \item \textbf{Validação temporal}, como no artigo:
    \begin{itemize}
      \item ajuste em 2017--2022;
      \item escolha de hiperparâmetros por otimização bayesiana (Optuna/TPE) \cite{akiba2019}, com \textit{early stopping} e recalibração isotônica em 2023;
      \item teste único, intocado, em 2024 (o ano mais recente com nascimentos completos).
    \end{itemize}
  \item \textbf{Desbalanceamento} (prevalência de $\approx$0,8\%): o peso de classe é tratado como hiperparâmetro. Não se usa sobreamostragem, para preservar a calibração.
\end{itemize}

\section{Avaliação}

\begin{itemize}
  \item \textbf{Métricas}, com IC 95\% por bootstrap estratificado:
    \begin{itemize}
      \item AUROC e AUPRC;
      \item F1, precisão, sensibilidade e especificidade no limiar escolhido na validação;
      \item escore de Brier e razão observado/esperado;
      \item fração dos óbitos entre os 5\% de maior risco, métrica comparável ao artigo.
    \end{itemize}
  \item \textbf{Comparações:} $\Delta$AUROC e $\Delta$AUPRC entre (A) e (B) e entre XGBoost e regressão logística, por bootstrap pareado.
  \item \textbf{Interpretação e subgrupos:} valores SHAP \cite{lundberg2017}, comparados ao ranking do artigo, e curvas de calibração por subgrupo.
  \item \textbf{Sensibilidade ao linkage:} repetição do modelo usando só as duas chaves mais restritivas.
\end{itemize}

\section{Entregas e cronograma}

\begin{center}
\begin{tabular}{ll}
\toprule
Etapa & Situação \\
\midrule
Download, linkage e base analítica & concluído \\
Modelos, ajuste de hiperparâmetros e avaliação & concluído \\
README (Parte 1) e notebook comentado (Parte 2) & concluído \\
Relatório no formato SBC ($\leq$ 12 páginas) & concluído \\
Apresentação (15 min + 5 de perguntas) & slides prontos \\
\bottomrule
\end{tabular}
\end{center}

\textbf{Limitações previstas.}
\begin{itemize}
  \item Cerca de 10\% dos óbitos não são pareados. Eles ficam rotulados como sobreviventes, o que tende a subestimar o desempenho.
  \item Os falsos pares introduzem ruído no rótulo.
  \item O modelo pré-parto usa a idade gestacional do parto, que é uma aproximação da estimativa disponível na admissão.
\end{itemize}

\bibliographystyle{sbc}
\bibliography{refs}

\end{document}
